\scanpage{13}
\mspage{6}
Revenons au cas de la catégorie isotopique
des surfaces orientées \deleted{de t.f.} \inserted{anabéliennes}, il lui est donc
associée une catégorie ``\unsure{complétée} profinie'',
avec des $\operatorname{Hom}^{\wedge}$ donnés par (12) --- ce sont
\deleted{(à $\unsure{\Pi_0|\pi_0}$ \unsure{compts} \unsure{dénombrables})}
des espaces sommes top. d'espaces compacts
totalement discontinus --- en fait, ce sont
des espaces loc. compacts dénombrables à l'infini,
totalement disconnexes, métrisables --- toutes
les qualités ! (Et pas de pts isolés, par
surcroît\ldots)

Limitons-nous maintenant au cas des
surfaces $X$ qui sont munies d'une structure
de courbe alg. complexe, compatible avec sa
topologie. Alors les morphismes entre tels
$X$, dans un $\operatorname{Hom}^{\wedge}_{\mathrm{isot}}$, et leur composition,
définis ici en utilisant de façon essentielle
la topologie du corps des complexes $\mathbb C$
(permettant d'associer aux schémas en courbes
lisses quasi-projectives sur $\unsure{C|\mathbb C}$ des surfaces
topologiques orientées de t.f.), peuvent se
décrire de façon purement algébro-géomé-
trique, indépendamment de cette topologie
sur $\mathbb C$ (cas anabélien, tjs). Pour le
\margin{12}
